\chapter{Architecture générale de la plateforme}
\label{chap:archi}

\section{Vue d'ensemble}
PrognoSense adopte une architecture \textbf{client-serveur en trois couches}, conçue pour
la modularité et la séparation des responsabilités :
\begin{itemize}[leftmargin=1.5em]
  \item un \textbf{backend} Python reposant sur le framework \textbf{FastAPI}, qui expose
        une API REST et un canal temps réel WebSocket, et qui encapsule la logique métier
        (traitement du signal, modèles d'IA, persistance) ;
  \item un \textbf{frontend} \textbf{React} (bundler Vite) offrant une interface
        opérateur réactive, organisée en sept pages fonctionnelles ;
  \item une couche de \textbf{persistance} (SQLAlchemy sur SQLite, compatible
        PostgreSQL/TimescaleDB) et des journaux (audit JSONL, fichiers modèles).
\end{itemize}

\begin{figure}[h]
\centering
\begin{tikzpicture}[
  font=\footnotesize,
  layer/.style={rectangle, rounded corners=4pt, draw=psbleu, fill=psbleu!7,
                minimum width=14cm, minimum height=1.5cm, align=center},
  comp/.style={rectangle, rounded corners=3pt, draw=psgris!60, fill=white,
               minimum height=0.85cm, align=center, font=\scriptsize},
  arr/.style={{Stealth[length=2mm]}-{Stealth[length=2mm]}, psgris, thick}]

  \node[layer, fill=psvert!8, draw=psvert] (front) {};
  \node[anchor=west, font=\bfseries\scriptsize, text=psvert] at (front.north west) {FRONTEND (React / Vite)};
  \node[comp] at (-5.2,0) {Vue Globale};
  \node[comp] at (-3.3,0) {Monitoring};
  \node[comp] at (-1.5,0) {IA Lab};
  \node[comp] at (0.3,0) {Pronostic};
  \node[comp] at (2.1,0) {Maintenance};
  \node[comp] at (4.1,0) {Audit};
  \node[comp] at (5.7,0) {Config};

  \node[layer, below=1.1cm of front] (back) {};
  \node[anchor=west, font=\bfseries\scriptsize, text=psbleu] at (back.north west) {BACKEND (FastAPI --- $\sim$20 routeurs REST + WebSocket)};
  \node[comp] at ($(back.center)+(-5.4,-0.05)$) {predict / model};
  \node[comp] at ($(back.center)+(-3.2,-0.05)$) {simulation (WS)};
  \node[comp] at ($(back.center)+(-1.1,-0.05)$) {spectral / ISO};
  \node[comp] at ($(back.center)+(1.0,-0.05)$) {ingest / workorders};
  \node[comp] at ($(back.center)+(3.3,-0.05)$) {kpi / analytics};
  \node[comp] at ($(back.center)+(5.3,-0.05)$) {copilot / auth};

  \node[layer, below=1.1cm of back, fill=psambre!10, draw=psambre] (ml) {};
  \node[anchor=west, font=\bfseries\scriptsize, text=psambre] at (ml.north west) {MODULES ML / SIGNAL};
  \node[comp] at ($(ml.center)+(-5.0,-0.05)$) {Model Registry};
  \node[comp] at ($(ml.center)+(-2.8,-0.05)$) {Feature Extractor};
  \node[comp] at ($(ml.center)+(-0.6,-0.05)$) {Autoencoder};
  \node[comp] at ($(ml.center)+(1.5,-0.05)$) {Anomaly Ensemble};
  \node[comp] at ($(ml.center)+(3.9,-0.05)$) {Drift / RAG};

  \node[layer, below=1.1cm of ml, fill=psgris!8, draw=psgris] (db) {};
  \node[anchor=west, font=\bfseries\scriptsize, text=psgris] at (db.north west) {PERSISTANCE};
  \node[comp] at ($(db.center)+(-4.0,-0.05)$) {SQLAlchemy / SQLite (WAL)};
  \node[comp] at ($(db.center)+(0.5,-0.05)$) {Audit JSONL};
  \node[comp] at ($(db.center)+(4.0,-0.05)$) {Modèles (.pkl/.pt)};

  \draw[arr] (front) -- (back) node[midway, right, font=\scriptsize] {REST / WebSocket};
  \draw[arr] (back) -- (ml);
  \draw[arr] (ml) -- (db);
\end{tikzpicture}
\caption{Architecture en couches de PrognoSense.}
\label{fig:archi}
\end{figure}

\section{Principes de conception}
Trois principes structurent l'implémentation :
\begin{description}[leftmargin=1.5em, style=nextline]
  \item[Registre de modèles centralisé.] Un singleton \texttt{ModelRegistry} charge tous
        les modèles en mémoire au démarrage, gère le \emph{modèle actif} par dataset
        (choix persisté), et expose une interface de prédiction unifiée. C'est le point de
        passage unique de toute inférence.
  \item[Gestionnaire de flotte.] Un singleton \texttt{FleetHealthManager} maintient, pour
        chaque machine, un historique circulaire (indice de santé, RUL, anomalie) et calcule
        tendances et alertes en temps réel.
  \item[Pipeline unifié.] L'adaptateur \texttt{UnifiedSignal} (chapitre~\ref{chap:datasets})
        garantit que la même chaîne de traitement s'applique à toutes les natures de
        signaux.
\end{description}

\section{Communication temps réel}
La surveillance en direct repose sur un \textbf{WebSocket} (\texttt{/api/ws/simulation})
plutôt que sur du \emph{polling}, afin de garantir une faible latence sans surcharge. Un
\texttt{ConnectionManager} gère plusieurs connexions simultanées, chacune avec son propre
état (machine, dataset, vitesse). À chaque cycle, le serveur calcule l'indice de santé,
génère un signal, exécute les prédictions et diffuse un message structuré au client. Le
client implémente une \textbf{reconnexion automatique à back-off exponentiel} (1, 2, 4,
\dots, 30\,s) pour la robustesse.

\section{Technologies retenues et justification}
\begin{table}[h]
\centering\small
\caption{Pile technologique et justification des choix.}
\label{tab:stack}
\begin{tabular}{>{\bfseries}p{3.0cm} p{6.0cm} p{4.6cm}}
\toprule
Brique & Technologie & Justification \\
\midrule
API backend & FastAPI + Uvicorn & Asynchrone, typage Pydantic, doc Swagger automatique \\
Traitement signal & NumPy, SciPy & FFT, Hilbert, filtres, intégration --- référence scientifique \\
ML classique & scikit-learn, XGBoost & Modèles éprouvés, rapides à entraîner et interpréter \\
ML profond & PyTorch & Auto-encodeur et CNN 1D \\
Persistance & SQLAlchemy + SQLite & Zéro infrastructure, portable vers PostgreSQL/Timescale \\
Connectivité & asyncua, paho-mqtt & Standards industriels OPC-UA et MQTT \\
IA générative & API Mistral + TF-IDF & Copilot et recherche documentaire (RAG) \\
Sécurité & python-jose, passlib & JWT et hachage bcrypt \\
Frontend & React, Vite, Recharts & Interface réactive, graphiques temps réel \\
\bottomrule
\end{tabular}
\end{table}

Le choix de \textbf{FastAPI} est motivé par sa nature asynchrone (essentielle pour le
WebSocket temps réel), sa validation de données intégrée (Pydantic) qui sécurise les
entrées, et sa documentation interactive automatique. Le choix de \textbf{SQLite} (plutôt
qu'une base lourde) relève d'une décision assumée : zéro dépendance d'infrastructure pour le
prototype, avec une bascule vers PostgreSQL/TimescaleDB possible sans modification du code
(chapitre~\ref{chap:indus}).
